\documentclass[UTF8,a4paper,11pt,openany]{ctexbook}
\usepackage{../common/statlearnbook}
\SLSetBookMetadata{第 4 册}{统计学习方法讲义 v2.6.0}{无监督学习与矩阵分解}
\begin{document}
\setcounter{page}{595}
\setcounter{chapter}{26}
\setcounter{figure}{0}
\pagestyle{headings}
\chaptermark{奇异值分解}
\SLAdvancedSection{代数推导与几何解释}{sec:V4-C03-S06}
\begin{geometrybox}
\textbf{几何解释。}单位球先经$V^T$旋转或反射，仍为单位球；再经$\Sigma$沿正交坐标轴缩放为半轴长度$\sigma_1,\ldots,\sigma_p$的椭球；最后经$U$旋转或反射到输出空间。
\end{geometrybox}
如\cref{fig:V4-C03-svd-geometry}所示，两个正交变换只改变坐标方向，所有长度伸缩都集中在$\Sigma$的非负对角元上。
% v2.3.1 figure UID: FIG-P467-01
% Proposed post-v2.3.1 polish for FIG-P467-01: protect key-label size and stroke hierarchy.
\tikzset{slfig-FIG-P467-01/.style={font=\fontsize{9.2pt}{11.0pt}\selectfont,every path/.append style={line cap=round,line join=round}}}
\begin{figure}[H]
\centering
\begin{tikzpicture}[slfig-FIG-P467-01,
  curve trace/.style={draw=SLGray!80,line width=.52pt},
  principal vector/.style={-{Stealth[length=1.95mm,width=1.18mm]},draw=SLBlue,line width=.98pt},
  gray basis/.style={-{Stealth[length=1.72mm,width=1.05mm]},draw=SLGray!88!black,line width=.68pt},
  teal basis/.style={-{Stealth[length=1.72mm,width=1.05mm]},draw=SLTeal,line width=.68pt},
  phase arrow/.style={-{Stealth[length=2.1mm,width=1.25mm]},draw=SLGold,line width=.86pt},
]
\begin{groupplot}[
  group style={group size=4 by 1,horizontal sep=8mm},
  width=.20\linewidth,height=.20\linewidth,
  xmin=-1.9,xmax=1.9,ymin=-1.9,ymax=1.9,
  axis equal image,axis lines=middle,axis line style={draw=SLInk,line width=.58pt},
  xtick=\empty,ytick=\empty,clip=false,
  title style={font=\bfseries\fontsize{9.4pt}{11.2pt}\selectfont,yshift=1mm},
]
\nextgroupplot[title={输入：单位圆}]
  \addplot[curve trace,domain=0:360,samples=181] ({cos(x)},{sin(x)});
  \addplot[principal vector] coordinates {(0,0) (.8192,.5736)};
  \addplot[gray basis] coordinates {(0,0) (1,0)};
  \addplot[teal basis] coordinates {(0,0) (0,1)};
  \addplot[only marks,mark=*,mark size=1.35pt,draw=SLBlue!58,fill=SLBlue!58]
    coordinates {(.5,.866) (-.5,.866) (-.866,-.5) (.866,-.5)};
\nextgroupplot[title={$V^{\mathsf T}$：正交旋转}]
  \addplot[curve trace,domain=0:360,samples=181] ({cos(x)},{sin(x)});
  \addplot[principal vector] coordinates {(0,0) (.9962,.0872)};
  \addplot[gray basis] coordinates {(0,0) (.8660,-.5000)};
  \addplot[teal basis] coordinates {(0,0) (.5000,.8660)};
\nextgroupplot[title={$\Sigma$：轴向缩放}]
  \addplot[curve trace,domain=0:360,samples=181] ({1.6*cos(x)},{.65*sin(x)});
  \addplot[principal vector] coordinates {(0,0) (1.5939,.0567)};
  \addplot[gray basis] coordinates {(0,0) (1.3856,-.3250)};
  \addplot[teal basis] coordinates {(0,0) (.8000,.5629)};
\nextgroupplot[title={$U$：正交旋转}]
  \addplot[curve trace,domain=0:360,samples=181]
    ({1.4501*cos(x)-.2747*sin(x)},{.6762*cos(x)+.5891*sin(x)});
  \addplot[principal vector] coordinates {(0,0) (1.4205,.7248)};
  \addplot[gray basis] coordinates {(0,0) (1.3931,.2910)};
  \addplot[teal basis] coordinates {(0,0) (.4871,.8483)};
\end{groupplot}
\draw[phase arrow] (group c1r1.east)--(group c2r1.west);
\draw[phase arrow] (group c2r1.east)--(group c3r1.west);
\draw[phase arrow] (group c3r1.east)--(group c4r1.west);
\node[font=\fontsize{9pt}{10.6pt}\selectfont,text=SLGray!94!black] at ([yshift=-6mm]group c2r1.south east) {蓝向量与两条基方向贯穿四个面板；坐标范围一致};
\end{tikzpicture}
\caption{SVD把线性变换分成正交变换、轴向缩放和正交变换}\label{fig:V4-C03-svd-geometry}
\end{figure}

\end{document}
